% ============================================================
% main.tex — CDS534 小组项目报告 · 封面+目录+合章（仅 E 可改）
% Topic 5 · Ride-Hailing Order & Trip Dispatch（网约车订单与派单数据库）
% 编译：Overleaf pdfLaTeX，以无错为准（warning 可留）
% 硬性要求：正文 10–12 页 · 字号 12 · 封面列组名与全员学号 · ER 矢量
% 写作约定：05_report/README.md + 01_docs/glossary.md（术语表）
% 章节页数预算：§1 1.5 / §2 2.0 / §3 1.5 / §4 3.0 / §5 1.5 / §6 0.5–1.0
% 超页压缩顺序：§4 问答表述 → 截图并排 → ER 图与目标清单不压缩
% ============================================================
\documentclass[12pt,a4paper]{article}

\usepackage[utf8]{inputenc}
\usepackage[T1]{fontenc}
\usepackage[a4paper,margin=2.5cm]{geometry}
\usepackage{graphicx}
\usepackage{booktabs}      % §5 结果表统一三线表
\usepackage{listings}      % SQL/脚本排版
\usepackage{xcolor}
\usepackage{url}
\usepackage[colorlinks=true,linkcolor=blue,urlcolor=blue,citecolor=blue]{hyperref}

% ---- 代码样式：等宽 9 号（组内约定）----
\lstset{
  basicstyle=\ttfamily\fontsize{9}{11}\selectfont,
  keywordstyle=\color{blue},
  commentstyle=\color{gray},
  numbers=left,
  numberstyle=\tiny\color{gray},
  breaklines=true,
  frame=single,
  language=SQL,
  showstringspaces=false,
  tabsize=2
}

% ---- 图片：截图统一裁边、同比例缩放；源文件在 03_data/screenshots/ ----
\graphicspath{{figures/}}

\begin{document}

% ================= 封面 =================
\begin{titlepage}
  \centering
  \vspace*{2cm}
  {\LARGE\bfseries Ride-Hailing Order and Trip Dispatch Database\\[0.3em]
   \large (Topic 5)}\par
  \vspace{1.5cm}
  {\large CDS 534 Database Management --- Group Project}\par
  \vspace{2.5cm}
  % TODO(胡中远): 组名（kickoff 遗留项，群内确认后回填；文件名 CDS534 Group TeamName.pdf 同步）
  {\large Group: \textbf{TODO-GroupName}}\par
  \vspace{1cm}
  % TODO(胡中远): 全员学号（顺序按角色 A–E，与附录 A 贡献表一致）
  \begin{tabular}{ll}
    \textbf{A \middot Business \& Requirements} & TODO 陈祎楠 (SID) \\
    \textbf{B \middot ER Modelling}             & TODO 张奕杰 (SID) \\
    \textbf{C \middot Implementation}           & TODO 陈妍吾 (SID) \\
    \textbf{D \middot Queries \& Experiments}   & TODO 凌彤 (SID) \\
    \textbf{E \middot PM \& Integration}        & TODO 胡中远 (SID) \\
  \end{tabular}\par
  \vfill
  {\large \today}\par
\end{titlepage}

\tableofcontents
\newpage

% ================= 正文（各人只写自己的章节文件）=================
% ============================================================
% §1 Business Context and Goals — A 主笔、E 润色 · 页数预算 1.5
% 素材：A1 业务流程 / A2 利益相关者 / A3 目标与约束 / A4 数据集选型 / A6 CAP 权衡
% 对应 rubric：Problem framing & significance 4 分 + Methodology 3 分（目标必须带具体数字）
% ============================================================
\section{Business Context and Goals}
\label{sec:context}

\subsection{Background and Business Process}
% TODO(A): 网约车运作链路（下单→派单→行程→支付→评价），配 A1 业务流程图（draw.io 导出矢量 PDF）
% 口径提醒：匹配队列不落库（D-05）——"下单量"分析属已知局限，§6 呼应
% 图：\begin{figure}[ht] \centering \includegraphics[width=0.9\textwidth]{figures/business_flow.pdf}
%      \caption{Business process of order dispatch and fulfilment} \label{fig:business-flow} \end{figure}
TODO(A): 段落。

\subsection{Stakeholders and Their Data Needs}
% TODO(A): A2 五方×2 条（乘客/司机/运营/财务/客服）压缩为一张 booktabs 表，标注每条需求对应的数据来源表
TODO(A): 表~\ref{tab:stakeholders}。

\subsection{Quantified Goals, Constraints and Trade-offs}
% TODO(A): A3 的 G1–G5 / C1–C6 摘要（P95<200ms、50 TPS、180 天保留、Trip ≥ 50,000 等；数字以全组表决定稿为准）
% TODO(A): A6 一致性/可用性/延迟权衡——派单与支付为何强一致（转账类比），报表可弱一致（读库/缓存）
TODO(A): 段落。

\subsection{Dataset Selection Rationale}
% TODO(A): A4 结论——生成器路线（faker + 固定 seed，分布参照 NYC TLC），样本规模见 A4 第三节
% 引用 \cite{tlctripdata}；TLC 字段适配一票否决论证压缩为 2–3 句（缺乘客身份/司机-车辆 1:1 源字段）
TODO(A): 段落。
    % §1 A 主笔、E 润色
% ============================================================
% §2 Entity-Relationship Model — B 主笔 · 页数预算 2.0 · 素材 B2–B5
% ER 关（11-02）检查项：官方约定遵循 / 每个扩展实体有业务理由 / 与 §1 目标呼应
% ============================================================
\section{Entity-Relationship Model}
\label{sec:er}

\subsection{Notation and Conventions}
% TODO(B): 官方约定遵循说明——矩形实体、属性入框、主键斜体、外键以 *id 标注、连线标基数、N:M 拆关联实体
TODO(B): 段落。

\subsection{Baseline Model: Five Entities}
% TODO(B): 五实体 ER 矢量图（draw.io 导出 PDF，勿用位图）+ 实体/属性/键说明
% 关键基数：Passenger–Trip 1:N / Driver–Trip 1:N / Driver–Vehicle N:1（vehicle_id UNIQUE）/ Trip–Payment 1:1（trip_id UNIQUE）
TODO(B): 图~\ref{fig:er-baseline}。

\subsection{Extensions}
% TODO(B): B3 表决通过的扩展，每张扩展表写一句业务理由 + 对查询的价值；未通过的方案写一句取舍说明
% 提示：扩展表 DDL 一律进 build.sql 第 3 节扩展区，底座永不修改（D-07/协作铁律）
TODO(B): 段落。

\subsection{Normalization}
% TODO(B): B5 逐表 1NF/2NF/3NF 判定与理由 + 一个传递依赖消除示例（用本项目表举例，供概念题 Q2 复用）
TODO(B): 段落。
         % §2 B 主笔
% ============================================================
% §3 Implementation — C 主笔 · 页数预算 1.5 · 素材 C1/C3/C4/C5
% SQL 关（11-02）检查项：干净环境从零复跑 build.sql + 导数成功
% ============================================================
\section{Implementation}
\label{sec:impl}

\subsection{Environment and Architecture}
% TODO(C): 环境与版本登记（env_record：最低基线 8.0.45）、GitHub 协作模式（脚本唯一事实、数据同源 seed）
TODO(C): 段落。

\subsection{DDL Organization and Key Constraints}
% TODO(C): build.sql 四节结构（建库 / 官方底座 / 扩展区 / 自检）；底座铁律；
%          UNIQUE 强制 1:1 与 N:1（trip_id UNIQUE / vehicle_id UNIQUE）；CHECK amount>=0 的版本兼容说明
TODO(C): 段落。

\subsection{Data Generation and Import}
% TODO(C): import_data.py 固定 seed=42，可重复执行；
%          行数：Passenger 3,000 / Driver 150 / Vehicle 150 / Trip 60,000 / Payment ≈ 51,300（A4 第三节参数）
TODO(C): 段落。

\subsection{Verification}
% TODO(C): 干净复跑记录 + 全套截图（行数自检、SHOW CREATE TABLE、样例查询），按 C2 截图规范命名归档
TODO(C): 段落。
       % §3 C 主笔
% ============================================================
% §4 SQL Queries and Conceptual Q&A — D 主笔 · 页数预算 3.0 · 素材 D2/D3/D4
% SQL 关检查项：报告中每条 SQL 与 02_sql/queries/ 脚本一字不差；截图可追溯
% ============================================================
\section{SQL Queries and Conceptual Q\&A}
\label{sec:sql}

\subsection{Business Queries}
% TODO(D): 查询清单 v1 共 8 条（官方 Q1–Q3 保底 + 进阶 Q4–Q8）；若 B3"上车区域"扩展通过，追加 q09 热区查询
% 每条固定四件套：业务问题（1 句）/ SQL（lstlisting，与脚本一致）/ 结果截图 / 业务解读 2–3 句（讲含义，不复述数字）
% 解读口径提醒：Q2 是"完成未支付"的对账催收，不含 canceled（A1 说明）；取消率按已入库行程计算（D-05 局限）
TODO(D): 逐条展开。

\subsection{Conceptual Q\&A}
% TODO(D): 8 道概念问答，每题 150–250 词，用本项目 schema 举例（不照抄官方答案表述）
% A/B/C 各复核 2 题；Q5 丢失更新记得写清 read-modify-write 前提（A 复核意见）
TODO(D): 逐题展开。
        % §4 D 主笔
% ============================================================
% §5 Evaluation — D 主笔、C 供数据 · 页数预算 1.5 · 素材 A5/D5/C6
% 数字关（11-02）检查项：结论句挂的数字与截图/日志逐一对账
% 实验数字统一由 D 本机产出（协作纪律 4），避免跨机性能差异
% ============================================================
\section{Evaluation}
\label{sec:eval}

\subsection{Workload Profile}
% TODO(D): A5 负载画像——Top5 高频查询（Q1 对账/Q6 调度/Q3、Q7 报表/Q5 催收）与读写比
TODO(D): 表。

\subsection{Experiment 1: Index Before/After}
% TODO(D): 选大表查询，EXPLAIN + 耗时成对截图；对照 G2 靶子（P95 < 200 ms）
% 候选索引：Trip(status)、Trip(start_time)、扩展通过时 Trip/关联表(zone_id)
TODO(D): 结果表 + 结论。

\subsection{Experiment 2: JOIN vs Subquery}
% TODO(D): 同一业务问题两种写法的执行计划对比（如 Q2 NOT EXISTS vs LEFT JOIN ... IS NULL）
TODO(D): 结果表 + 结论。

\subsection{Experiment 3: Concurrent Transactions}
% TODO(D): 两会话模拟并发支付，展示丢失更新与隔离级别效果；按 G3（50 TPS 设计余量）加压
TODO(D): 结果表 + 结论。

\subsection{Results Summary and Goal Recheck}
% TODO(D): 统一结果表样收尾；每条结论必须挂实验数字，§6 与 §1 目标闭环
TODO(D): 段落。
       % §5 D 主笔、C 供数据
% ============================================================
% §6 Conclusion and Limitations — A+E 合写 · 页数预算 0.5–1.0
% rubric：Evaluation 结论有据 3 分——与 §1 目标逐条闭环
% ============================================================
\section{Conclusion and Limitations}
\label{sec:conclusion}

\subsection{Goal Recheck}
% TODO(A/E): 对照 §1 的 G1–G5 逐条回收：达成/未达成 + 挂实验编号
TODO(A+E): 段落。

\subsection{Limitations}
% TODO(A/E): 已知局限——①仿真数据：跨组对比结论是生成规则回显，不做因果声称（D-10 撰写提醒）
%            ②D-05：匹配期取消不落库，"下单量/应答率"不可分析 ③单机环境性能数字不可外推 ④样本规模
TODO(A+E): 段落。

\subsection{Future Work}
% TODO(A/E): 分区表 / 读写分离 / NoSQL 对比（方案给定方向）
TODO(A+E): 段落。
 % §6 A+E 合写

% ================= References =================
% TODO(全员): 正文落笔时用 \cite{...} 替换 \nocite{*}（骨架期先全量展示引用库 v1）
\nocite{*}
\bibliographystyle{plain}
\bibliography{references}

% ================= 附录 =================
\appendix
% ============================================================
% 附录 A 贡献表 — E 制表、全员确认（11-03 定稿）
% 内容依据：01_docs 各任务交付记录；工作量% 五人合计 100%
% ============================================================
\section{Contribution Statement}
\label{app:contribution}

% TODO(胡中远): 11-03 前按仓库提交记录与任务表填写，全员逐一确认后冻结
\begin{table}[ht]
\centering
\begin{tabular}{llll}
\toprule
Member (SID) & Role & Main Contributions & Load \\
\midrule
CHEN Yinan (3162550)  & A & TODO & TODO\% \\
ZHANG Yijie (3162641) & B & TODO & TODO\% \\
CHEN Yanwu (3162744)  & C & TODO & TODO\% \\
LING Tong (3162677)   & D & TODO & TODO\% \\
HU Zhongyuan (3162304)& E & TODO & TODO\% \\
\bottomrule
\end{tabular}
\caption{Team contributions}
\label{tab:contribution}
\end{table}
   % 附录 A 贡献表（E）
% ============================================================
% 附录 B 完整脚本 — C 供内容、E 排版（页数上限可放宽，不计 10–12 页正文口径）
% ============================================================
\section{Full Build Script}
\label{app:sql}

% TODO(C): 定稿时把 build.sql 全文复制为 05_report/figures/appendix_build.sql（或同目录），
% 取消下行注释即可整文件排版；导数脚本只给入口说明（python import_data.py + seed）+ 关键参数注释
% \lstinputlisting[language=SQL,caption={build.sql --- full DDL with repeatability self-check}]{figures/appendix_build.sql}

TODO(C): 待 11-03 前提供 build.sql 定稿全文与 import\_data.py 入口说明。
            % 附录 B 完整脚本（C 供、E 排版）

\end{document}
